% ============================================================================
% Captions for the two rebuilt resource figures.  Every number below is taken
% from p2_figs/figure_numbers.json, which make_figs.py writes only after all
% 36 self-verification checks pass.  Nothing here is typed by hand.
% ============================================================================

% ---------------------------------------------------------------- FIGURE 1
\begin{figure*}[tp]\centering
% ===== fig:decoupling (REBUILT 2026-09-18) -- native pgfplots, no external .dat, no hand-typed numbers.
% Emitted by p2_figs/make_figs.py from data/cost_vs_ent.json, data/n19_suite.json,
% data/resource_master.json and p2_calc/stats_resource.json; all statistics re-verified at emit time.
\definecolor{dcCov}{HTML}{E8543A}\definecolor{dcMrf}{HTML}{10A87B}\definecolor{dcIon}{HTML}{2E6FE0}
\definecolor{dcTrend}{HTML}{9A9A9A}\definecolor{dcPool}{HTML}{6E6E6E}\definecolor{inkPrim}{HTML}{1B1B1B}
\definecolor{scAh}{HTML}{D55E00}\definecolor{scAd}{HTML}{F0A868}\definecolor{scBh}{HTML}{00755E}\definecolor{scBd}{HTML}{5FC2A5}\definecolor{scCh}{HTML}{1B4F9C}\definecolor{scCd}{HTML}{79AEE8}
\begin{tikzpicture}
\begin{groupplot}[group style={group size=2 by 2, horizontal sep=1.8cm, vertical sep=2.05cm},
  width=7.45cm, height=5.6cm, paperaxis,
  title style={at={(0,1)},anchor=south west,font=\small\bfseries,yshift=2pt},
  label style={font=\normalsize}, tick label style={font=\footnotesize},
  ymajorgrids, grid style={draw=black!8}]
% ---- (a) molecules: F1 and |S| rise together (the regime where F1 looks predictive) ----
\nextgroupplot[title={(a)~molecules: they rise together},
  xlabel={$\mathcal F_1=4\,\mathrm{tr}[\gamma(1{-}\gamma)]$}, ylabel={determinant support $|\mathcal S|$},
  ymode=log, xmin=-0.9, xmax=12.29, ymin=0.62, ymax=460, enlarge x limits=false]
\addplot[dcTrend,line width=1pt,dashed,domain=0.000:11.377,samples=60,forget plot] {2^(0.4736*x+0.8502)};
\addplot[only marks,mark=*,mark size=1.9pt,draw=white,line width=0.4pt,fill=dcCov] coordinates {(0.0380,1) (2.5114,2) (0.0008,1) (3.4603,4) (0.0676,1) (3.9849,2) (0.6882,8) (8.1924,46) (0.8474,7) (11.3765,80) (0.0392,1) (3.7542,6) (0.0068,1) (6.9358,13) (0.0803,1) (10.2890,68) (0.2515,2) (5.1131,11) (0.5214,9) (7.7351,91)};
\addplot[only marks,mark=*,mark size=1.9pt,draw=white,line width=0.4pt,fill=dcMrf] coordinates {(4.7428,13) (8.0009,6) (0.0127,1) (3.7026,5) (1.6728,11) (6.0001,4) (1.3066,27) (5.9378,9) (0.7900,5) (7.4487,25)};
\addplot[only marks,mark=*,mark size=1.9pt,draw=white,line width=0.4pt,fill=dcIon] coordinates {(0.0021,1) (0.0750,2) (0.0000,1) (0.0007,1) (0.5181,5) (0.5731,3) (0.0353,1) (0.3864,2)};
\node[inkPrim,font=\scriptsize,anchor=north west,align=left] at (rel axis cs:0.02,0.97) {Pearson $r=0.80$ on $\log_2|\mathcal S|$\\Spearman $\rho=0.862$ ($n=38$)};
% ---- (b) the six (L,N) strata against F1 ----
\nextgroupplot[title={(b)~Hubbard: $\mathcal F_1$ runs backwards},
  xlabel={one-body magic $\mathcal F_1$}, ylabel={determinant support $|\mathcal S|$},
  ymode=log, xmin=-0.9, xmax=20.2, ymin=40, ymax=900000, enlarge x limits=false]
\addplot[scAh,mark=*,mark size=2.1pt,line width=1.0pt,mark options={fill=scAh,draw=scAh}] coordinates {(0.6258,330) (2.2316,312) (5.9132,260) (9.6047,147) (11.2997,68)};
\addplot[scAd,mark=o,mark size=2.1pt,line width=1.0pt,mark options={fill=scAd,draw=scAd}] coordinates {(0.3396,200) (1.1047,196) (2.8100,188) (4.9945,168) (6.5809,143)};
\addplot[scBh,mark=square*,mark size=2.0pt,line width=1.0pt,mark options={fill=scBh,draw=scBh}] coordinates {(0.8020,3395) (2.9184,3120) (7.8807,2176) (12.8163,880) (15.0697,327)};
\addplot[scBd,mark=square,mark size=2.0pt,line width=1.0pt,mark options={fill=scBd,draw=scBd}] coordinates {(0.4855,2342) (1.6348,2272) (4.3683,2041) (7.9600,1600) (10.4394,1091)};
\addplot[scCh,mark=triangle*,mark size=2.6pt,line width=1.0pt,mark options={fill=scCh,draw=scCh}] coordinates {(0.9724,35678) (3.6002,30975) (9.8568,17732) (16.0306,5084) (18.8404,1476)};
\addplot[scCd,mark=triangle,mark size=2.6pt,line width=1.0pt,mark options={fill=scCd,draw=scCd}] coordinates {(0.6329,26503) (2.1803,25235) (5.9969,21203) (11.0065,13832) (14.3023,7527)};
\node[inkPrim,font=\scriptsize,anchor=north west,align=left] at (rel axis cs:0.02,0.97) {$U=1\to16$ along each curve\\within each sweep: $\rho=-1.000$\\pooled ($n=30$): $\rho=-0.109$ (withdrawn)};
% ---- (c) the same six strata against chi ----
\nextgroupplot[title={(c)~Hubbard: $\chi$ runs forwards},
  xlabel={minimal bond dimension $\chi$},
  ylabel={determinant support $|\mathcal S|$},
  ymode=log, xmin=2.6, xmax=18.4, ymin=40, ymax=900000, enlarge x limits=false,
  legend style={at={(1.16,-0.46)},anchor=north,legend columns=6,draw=black!25,font=\footnotesize},
  legend cell align=left]
\addplot[scAh,mark=*,mark size=2.1pt,line width=1.0pt,mark options={fill=scAh,draw=scAh}] coordinates {(12,330) (10,312) (6,260) (5,147) (4,68)};
\addlegendentry{$(6{,}6)$}
\addplot[scAd,mark=o,mark size=2.1pt,line width=1.0pt,mark options={fill=scAd,draw=scAd}] coordinates {(11,200) (10,196) (9,188) (9,168) (9,143)};
\addlegendentry{$(6{,}4)$}
\addplot[scBh,mark=square*,mark size=2.0pt,line width=1.0pt,mark options={fill=scBh,draw=scBh}] coordinates {(11,3395) (11,3120) (8,2176) (8,880) (7,327)};
\addlegendentry{$(8{,}8)$}
\addplot[scBd,mark=square,mark size=2.0pt,line width=1.0pt,mark options={fill=scBd,draw=scBd}] coordinates {(16,2342) (16,2272) (15,2041) (14,1600) (13,1091)};
\addlegendentry{$(8{,}6)$}
\addplot[scCh,mark=triangle*,mark size=2.6pt,line width=1.0pt,mark options={fill=scCh,draw=scCh}] coordinates {(17,35678) (14,30975) (10,17732) (6,5084) (6,1476)};
\addlegendentry{$(10{,}10)$}
\addplot[scCd,mark=triangle,mark size=2.6pt,line width=1.0pt,mark options={fill=scCd,draw=scCd}] coordinates {(13,26503) (13,25235) (13,21203) (11,13832) (9,7527)};
\addlegendentry{$(10{,}8)$}
\node[inkPrim,font=\scriptsize,anchor=north west,align=left] at (rel axis cs:0.02,0.97) {$\chi$ falls as $U$ rises\\within each sweep: $\rho=0.894$ to $1.000$\\pooled ($n=30$): $\rho=+0.569$ (withdrawn)};
% ---- (d) Simpson panel ----
\nextgroupplot[title={(d)~stratified vs pooled},
  xlabel={stratum $(L,N)$}, ylabel={Spearman $\rho$ vs $|\mathcal S|$},
  xmin=0.45, xmax=6.55, ymin=-1.42, ymax=1.95, ytick={-1,-0.5,0,0.5,1},
  xtick={1,2,3,4,5,6}, xticklabels={$(6{,}6)$,$(6{,}4)$,$(8{,}8)$,$(8{,}6)$,$(10{,}10)$,$(10{,}8)$},
  x tick label style={font=\scriptsize,rotate=30,anchor=east}, ymajorgrids=false]
\addplot[dcPool,densely dashed,line width=0.9pt,forget plot] coordinates {(0.45,0.5692) (6.55,0.5692)};
\addplot[dcPool,densely dashed,line width=0.9pt,forget plot] coordinates {(0.45,-0.1088) (6.55,-0.1088)};
\addplot[only marks,mark=*,mark size=2.4pt,scAh,mark options={fill=scAh,draw=white,line width=0.4pt},forget plot] coordinates {(1,1.0000)};
\addplot[only marks,mark=square*,mark size=2.3pt,scAh,mark options={fill=scAh,draw=white,line width=0.4pt},forget plot] coordinates {(1,-1.0000)};
\addplot[only marks,mark=*,mark size=2.4pt,scAd,mark options={fill=scAd,draw=white,line width=0.4pt},forget plot] coordinates {(2,0.8944)};
\addplot[only marks,mark=square*,mark size=2.3pt,scAd,mark options={fill=scAd,draw=white,line width=0.4pt},forget plot] coordinates {(2,-1.0000)};
\addplot[only marks,mark=*,mark size=2.4pt,scBh,mark options={fill=scBh,draw=white,line width=0.4pt},forget plot] coordinates {(3,0.9487)};
\addplot[only marks,mark=square*,mark size=2.3pt,scBh,mark options={fill=scBh,draw=white,line width=0.4pt},forget plot] coordinates {(3,-1.0000)};
\addplot[only marks,mark=*,mark size=2.4pt,scBd,mark options={fill=scBd,draw=white,line width=0.4pt},forget plot] coordinates {(4,0.9747)};
\addplot[only marks,mark=square*,mark size=2.3pt,scBd,mark options={fill=scBd,draw=white,line width=0.4pt},forget plot] coordinates {(4,-1.0000)};
\addplot[only marks,mark=*,mark size=2.4pt,scCh,mark options={fill=scCh,draw=white,line width=0.4pt},forget plot] coordinates {(5,0.9747)};
\addplot[only marks,mark=square*,mark size=2.3pt,scCh,mark options={fill=scCh,draw=white,line width=0.4pt},forget plot] coordinates {(5,-1.0000)};
\addplot[only marks,mark=*,mark size=2.4pt,scCd,mark options={fill=scCd,draw=white,line width=0.4pt},forget plot] coordinates {(6,0.8944)};
\addplot[only marks,mark=square*,mark size=2.3pt,scCd,mark options={fill=scCd,draw=white,line width=0.4pt},forget plot] coordinates {(6,-1.0000)};
\node[inkPrim,font=\scriptsize,anchor=north west] at (rel axis cs:0.02,0.995) {$\bullet$\; $\rho(\chi,|\mathcal S|)$ per stratum};
% D12: this key used to sit at the bottom-left, corner-to-corner with the (6,6) datum.
\node[inkPrim,font=\scriptsize,anchor=north west] at (rel axis cs:0.02,0.90) {$\blacksquare$\; $\rho(\mathcal F_1,|\mathcal S|)$ per stratum};
\node[dcPool,font=\scriptsize,anchor=north west] at (rel axis cs:0.035,0.578) {pooled ($n=30$): $\rho=+0.569$ (withdrawn)};
\node[dcPool,font=\scriptsize,anchor=north west] at (rel axis cs:0.035,0.377) {pooled ($n=30$): $\rho=-0.109$ (withdrawn)};
\end{groupplot}
\end{tikzpicture}

\caption{\label{fig:decoupling}\textbf{Inside every stratum one-body magic $\mathcal F_1$ orders the
static determinant support $|\mathcal S|$ perfectly and with the sign inverted; what fails is the
\emph{pooled} coefficient, not $\mathcal F_1$.}
\textbf{(a)}~Nineteen molecular active spaces at equilibrium and at dissociation ($38$ points):
in the static-correlation regime $\mathcal F_1$ and $|\mathcal S|$ rise together
(Pearson $r=0.80$ on $\log_2|\mathcal S|$; Spearman $\rho=0.862$). The grey dashed line is the least-squares fit $\log_2|\mathcal S|=0.474\,\mathcal F_1+0.850$ over those same points, a guide to that correlation and not a model of the cost.
\textbf{(b)}~The Hubbard interaction sweep is not one cloud of $30$ points but six strata
$(L,N)$ of five on-site repulsions $U=1,2,4,8,16$; one curve per stratum. Inside every one of
the six, raising $U$ raises $\mathcal F_1$ and lowers $|\mathcal S|$, and does so with a perfect
rank reversal: $\rho(\mathcal F_1,|\mathcal S|)=-1.000$ in all six. Pooling across strata
collapses the coefficient to $\rho=-0.109$; that collapse is a Simpson reversal, not an absence
of signal.
\textbf{(c)}~The same six strata against the minimal bond dimension: $\rho(\chi,|\mathcal S|)$
runs from $+0.894$ to $+1.000$ within strata, with a mean of $0.948$, and does not change sign between
them. Its pooled coefficient, like $\mathcal F_1$'s, is a Simpson reversal and is not used: pooling six
deterministic five-point sweeps destroys within-stratum structure whichever way the result falls, and
it fell in our favour on one of the two.
\textbf{(d)}~The six within-stratum coefficients against the two pooled ones. Stratified, the
$\mathcal F_1$ evidence is one-sided and exact: Kendall $\tau=-1.000$ with $0$ concordant of
$60$ within-stratum pairs and no ties, one-sided exact permutation $p=3.3\times10^{-13}$ over
the $(5!)^6=2.99\times10^{12}$ within-stratum relabellings. The reading
``$\rho=-0.11$ ($n=30$, $p\approx0.56$, not significant)'' is the pooled number and should be
replaced by the stratified test.
Two limits of this evidence, which the stratified test does not remove. First, the five points
of a stratum are a deterministic monotone sweep in $U$, so the effective number of independent
observations is $6$, not $30$; the exact $p$ counts relabellings, not independent samples.
Second, this is a statement about rank direction only: inside a fixed sweep $\mathcal F_1$
orders the five costs \emph{perfectly}, so the claim is not that $\mathcal F_1$ is uninformative
about $|\mathcal S|$ but that its sign is not stable---inverted here, positive over the
molecular set of panel~(a)---which is exactly what an orbital invariant compared against a
basis-dependent cost must do.}
\end{figure*}

% ---------------------------------------------------------------- FIGURE 2
\begin{figure*}[tp]\centering
% ===== fig:master (REBUILT 2026-09-18) -- 74 exact ground states, native pgfplots.
% Emitted by p2_figs/make_figs.py from data/resource_master.json and p2_calc/stats_resource.json.
% No hand-typed numbers: every coordinate, coefficient and interval is read from those files.
\definecolor{rmMol}{HTML}{0072B2}\definecolor{rmHub}{HTML}{E69F00}
\definecolor{rmChn}{HTML}{009E73}\definecolor{rmLad}{HTML}{D55E00}\definecolor{rmInk}{HTML}{1B1B1B}
\definecolor{rmGrey}{HTML}{6E6E6E}
\begin{tikzpicture}
\begin{groupplot}[group style={group size=2 by 2, horizontal sep=1.8cm, vertical sep=2.05cm},
  width=7.45cm, height=5.6cm, paperaxis,
  title style={at={(0,1)},anchor=south west,font=\small\bfseries,yshift=2pt},
  label style={font=\normalsize}, tick label style={font=\footnotesize},
  legend style={font=\scriptsize,draw=black!25,fill=white,fill opacity=0.92,text opacity=1},
  legend cell align=left]
% ---- (a) |S| vs chi, with the eleven molecular floor points marked ----
\nextgroupplot[title={(a)~cost tracks $\chi$}, xmode=log, ymode=log,
  xlabel={minimal bond dimension $\chi$}, ylabel={determinant support $|\mathcal S|$},
  xmin=0.8,xmax=70,ymin=0.50,ymax=900000, ymajorgrids,xmajorgrids, grid style={draw=black!8},
  legend to name=fig2legend, legend columns=3, legend style={draw=black!25,/tikz/every even column/.append style={column sep=11pt}}]
\addplot[only marks,mark=*,mark size=1.9pt,rmMol,mark options={fill=rmMol,draw=white,line width=0.3pt}] coordinates {(2,2) (4,2) (4,4) (2,2) (2,5) (2,3) (14,8) (30,46) (7,7) (32,80) (11,13) (2,6) (8,6) (13,13) (42,68) (4,2) (8,11) (8,9) (36,91) (2,2) (4,5) (4,11) (2,4) (17,27) (7,9) (11,5) (7,25)};
\addlegendentry{molecules, $|\mathcal S|>1$ ($\times27$)}
\addplot[only marks,mark=o,mark size=2.2pt,rmMol,line width=0.8pt] coordinates {(1,1) (1,1) (1,1) (2,1) (1,1) (1,1) (1,1) (1,1) (4,1) (1,1) (1,1)};
\addlegendentry{molecules at the floor $|\mathcal S|=1$ ($\times11$)}
\addplot[only marks,mark=square*,mark size=1.9pt,rmHub,mark options={fill=rmHub,draw=white,line width=0.3pt}] coordinates {(12,330) (10,312) (6,260) (5,147) (4,68) (11,200) (10,196) (9,188) (9,168) (9,143) (11,3395) (11,3120) (8,2176) (8,880) (7,327) (16,2342) (16,2272) (15,2041) (14,1600) (13,1091) (17,35678) (14,30975) (10,17732) (6,5084) (6,1476) (13,26503) (13,25235) (13,21203) (11,13832) (9,7527)};
\addlegendentry{Hubbard sweep ($\times30$)}
\addplot[only marks,mark=triangle*,mark size=2.2pt,rmChn,mark options={fill=rmChn,draw=white,line width=0.3pt}] coordinates {(5,147) (8,880) (6,5084)};
\addlegendentry{chain ($\times3$)}
\addplot[only marks,mark=diamond*,mark size=2.2pt,rmLad,mark options={fill=rmLad,draw=white,line width=0.3pt}] coordinates {(13,224) (38,1479) (38,11375)};
\addlegendentry{2-leg ladder ($\times3$)}
\addplot[rmInk,densely dashed,line width=0.9pt,mark=none,forget plot] coordinates {(2,2) (2,4) (2,8) (2,16) (2,32)};
\node[rmInk,font=\scriptsize,anchor=west] at (axis cs:2.15,120000) {witness: $\chi{=}2,\;|\mathcal S|{=}2^{K}$};
\node[rmInk,font=\scriptsize,anchor=south east] at (rel axis cs:0.98,0.02) {pooled $\rho=+0.722$};
% ---- (b) |S| vs F1, same 74 points, same marking ----
\nextgroupplot[title={(b)~$\mathcal F_1$ does not}, ymode=log,
  xlabel={one-body magic $\mathcal F_1$}, ylabel={determinant support $|\mathcal S|$},
  xmin=-1.2,xmax=23,ymin=0.50,ymax=900000, ymajorgrids, grid style={draw=black!8}]
\addplot[only marks,mark=*,mark size=1.9pt,rmMol,mark options={fill=rmMol,draw=white,line width=0.3pt}] coordinates {(2.5114,2) (0.0750,2) (3.4603,4) (3.9849,2) (0.5181,5) (0.5731,3) (0.6882,8) (8.1924,46) (0.8474,7) (11.3765,80) (4.7428,13) (8.0009,6) (3.7542,6) (6.9358,13) (10.2890,68) (0.2515,2) (5.1131,11) (0.5214,9) (7.7351,91) (0.3864,2) (3.7026,5) (1.6728,11) (6.0001,4) (1.3066,27) (5.9378,9) (0.7900,5) (7.4487,25)};
\addplot[only marks,mark=o,mark size=2.2pt,rmMol,line width=0.8pt] coordinates {(0.0380,1) (0.0021,1) (0.0008,1) (0.0676,1) (0.0000,1) (0.0007,1) (0.0392,1) (0.0068,1) (0.0803,1) (0.0353,1) (0.0127,1)};
\addplot[only marks,mark=square*,mark size=1.9pt,rmHub,mark options={fill=rmHub,draw=white,line width=0.3pt}] coordinates {(0.6258,330) (2.2316,312) (5.9132,260) (9.6047,147) (11.2997,68) (0.3396,200) (1.1047,196) (2.8100,188) (4.9945,168) (6.5809,143) (0.8020,3395) (2.9184,3120) (7.8807,2176) (12.8163,880) (15.0697,327) (0.4855,2342) (1.6348,2272) (4.3683,2041) (7.9600,1600) (10.4394,1091) (0.9724,35678) (3.6002,30975) (9.8568,17732) (16.0306,5084) (18.8404,1476) (0.6329,26503) (2.1803,25235) (5.9969,21203) (11.0065,13832) (14.3023,7527)};
\addplot[only marks,mark=triangle*,mark size=2.2pt,rmChn,mark options={fill=rmChn,draw=white,line width=0.3pt}] coordinates {(9.6047,147) (12.8163,880) (16.0306,5084)};
\addplot[only marks,mark=diamond*,mark size=2.2pt,rmLad,mark options={fill=rmLad,draw=white,line width=0.3pt}] coordinates {(9.6948,224) (12.9107,1479) (16.0889,11375)};
\node[rmInk,font=\scriptsize,anchor=north west,align=left] at (rel axis cs:0.02,0.97) {pooled $\rho=+0.599$ over the same 74 rows\\Hubbard branch: $U\!\uparrow\;\Rightarrow\;\mathcal F_1\!\uparrow,\;|\mathcal S|\!\downarrow$};
% ---- (c) the four coefficients, the intervals the deposit supports, and the tie ceiling ----
\nextgroupplot[title={(c)~coefficients; intervals for $n{=}38$ only},
  xlabel={Spearman $\rho$ against $|\mathcal S|$}, ylabel={},
  xmin=-0.03,xmax=1.03,ymin=-1.55,ymax=5.45, xtick={0,0.25,0.5,0.75,1},
  ytick={1,2,3,4}, yticklabels={{$\mathcal F_1$, $n{=}74$},{$\chi$, $n{=}74$},{$\mathcal F_1$, $n{=}38$},{$\chi$, $n{=}38$}},
  y tick label style={font=\scriptsize}, ymajorgrids=false, xmajorgrids, grid style={draw=black!8}]
\addplot[rmGrey,densely dashed,line width=0.9pt,forget plot] coordinates {(0.9863,0.45) (0.9863,4.45)};
\addplot[rmMol,line width=1.0pt,forget plot] coordinates {(0.8051,4) (0.9589,4)};
\addplot[rmMol,line width=1.0pt,forget plot] coordinates {(0.8051,3.87) (0.8051,4.13)};
\addplot[rmMol,line width=1.0pt,forget plot] coordinates {(0.9589,3.87) (0.9589,4.13)};
\addplot[rmLad,line width=1.0pt,forget plot] coordinates {(0.7091,3) (0.9294,3)};
\addplot[rmLad,line width=1.0pt,forget plot] coordinates {(0.7091,2.87) (0.7091,3.13)};
\addplot[rmLad,line width=1.0pt,forget plot] coordinates {(0.9294,2.87) (0.9294,3.13)};
\addplot[only marks,mark=*,mark size=2.4pt,rmMol,mark options={fill=rmMol,draw=white,line width=0.4pt},forget plot] coordinates {(0.9011,4)};
\addplot[only marks,mark=*,mark size=2.4pt,rmLad,mark options={fill=rmLad,draw=white,line width=0.4pt},forget plot] coordinates {(0.8618,3)};
\addplot[only marks,mark=*,mark size=2.4pt,rmMol,mark options={fill=rmMol,draw=white,line width=0.4pt},forget plot] coordinates {(0.7224,2)};
\addplot[only marks,mark=*,mark size=2.4pt,rmLad,mark options={fill=rmLad,draw=white,line width=0.4pt},forget plot] coordinates {(0.5992,1)};
\node[rmGrey,font=\scriptsize,anchor=east] at (axis cs:0.9863,5.05) {tie ceiling $0.986$\;};
\node[rmInk,font=\scriptsize,anchor=west] at (axis cs:-0.01,-0.40) {$n{=}38$: $\Delta\rho=0.039$ ($t=0.94$, $p=0.35$)};
\node[rmInk,font=\scriptsize,anchor=west] at (axis cs:-0.01,-1.15) {$n{=}74$: $\Delta\rho=0.123$ ($t=1.50$, $p=0.14$)};
% ---- (d) the molecular coefficient in rank space: the eleven-point tied block ----
\nextgroupplot[title={(d)~the 38 molecules in rank space},
  xlabel={midrank of $|\mathcal S|$}, ylabel={midrank of the predictor},
  xmin=0,xmax=40,ymin=0,ymax=59, xmajorgrids, ymajorgrids, grid style={draw=black!8},
  ytick={0,10,20,30,40},
  legend style={at={(0.98,0.985)},anchor=north east,font=\scriptsize}]
\addplot[rmGrey,densely dashed,line width=0.8pt,forget plot] coordinates {(1,1) (38,38)};
\addplot[rmGrey,line width=0.7pt,forget plot] coordinates {(6.0,25.0) (6.0,41.5)};
\addplot[only marks,mark=*,mark size=1.9pt,rmMol,mark options={fill=rmMol,draw=white,line width=0.3pt}] coordinates {(6.0,5.0) (14.0,13.5) (6.0,5.0) (14.0,20.5) (6.0,5.0) (18.5,20.5) (6.0,13.5) (14.0,13.5) (6.0,5.0) (6.0,5.0) (21.0,13.5) (17.0,13.5) (26.0,33.0) (35.0,35.0) (25.0,25.0) (37.0,36.0) (31.5,30.5) (23.5,13.5) (6.0,5.0) (23.5,28.0) (6.0,5.0) (31.5,32.0) (6.0,20.5) (36.0,38.0) (14.0,20.5) (29.5,28.0) (27.5,28.0) (38.0,37.0) (6.0,5.0) (14.0,13.5) (6.0,5.0) (21.0,20.5) (29.5,20.5) (18.5,13.5) (34.0,34.0) (27.5,25.0) (21.0,30.5) (33.0,25.0)};
\addlegendentry{$\chi$, $\rho=0.901$}
\addplot[only marks,mark=square*,mark size=1.9pt,rmLad,mark options={fill=rmLad,draw=white,line width=0.3pt}] coordinates {(6.0,8.0) (14.0,23.0) (6.0,4.0) (14.0,11.0) (6.0,3.0) (18.5,24.0) (6.0,10.0) (14.0,27.0) (6.0,1.0) (6.0,2.0) (21.0,15.0) (17.0,17.0) (26.0,18.0) (35.0,36.0) (25.0,20.0) (37.0,38.0) (31.5,28.0) (23.5,35.0) (6.0,9.0) (23.5,26.0) (6.0,5.0) (31.5,32.0) (6.0,12.0) (36.0,37.0) (14.0,13.0) (29.5,29.0) (27.5,16.0) (38.0,34.0) (6.0,7.0) (14.0,14.0) (6.0,6.0) (21.0,25.0) (29.5,22.0) (18.5,31.0) (34.0,21.0) (27.5,30.0) (21.0,19.0) (33.0,33.0)};
\addlegendentry{$\mathcal F_1$, $\rho=0.862$}
\node[rmInk,font=\scriptsize,anchor=north west,align=left] at (rel axis cs:0.02,0.985) {11 of 38 tied at $|\mathcal S|=1$\\midrank $6.0$: $|\rho|\le 0.986$};
% ---- D13: the grey rule marks the block of tied midranks; it now carries a label ----
\node[rmGrey,font=\scriptsize,anchor=west] at (axis cs:7.1,33.0) {tied block};
\end{groupplot}
\node[anchor=north,yshift=-1.25cm] at (group c1r2.south east) {\ref*{fig2legend}};
\end{tikzpicture}

\caption{\label{fig:master}\textbf{The resource thesis over the $74$ deposited rows
($71$ distinct exact ground states).} Nineteen molecules at equilibrium and dissociation, a
$30$-point Hubbard interaction sweep, and a matched chain/two-leg-ladder pair; all by exact
diagonalization.
\textbf{(a)}~The determinant support tracks the minimal bond dimension in every family (pooled
Spearman $\rho=+0.722$ over all $74$ rows, $+0.746$ over the $71$ distinct states---see
\emph{Provenance} below), with the geminal witness $\chi{=}2$, $|\mathcal S|{=}2^{K}$ the
deliberately loose lower bound. Open circles are the $11$ of $38$ molecular points pinned at the
floor $|\mathcal S|=1$---$29\%$ of the molecular set---where the support cannot resolve anything
and from which both molecular coefficients draw much of their leverage.
\textbf{(b)}~One-body magic does not ($\rho=+0.599$ over the same rows, $+0.592$ over the $71$
distinct states): the Hubbard branch runs
the wrong way ($U\!\uparrow$ raises $\mathcal F_1$ and lowers $|\mathcal S|$) and the
chain/ladder pair sits at $\mathcal F_1$ values agreeing to within $0.9\%$ while $|\mathcal S|$
differs by $1.52$, $1.68$ and $2.24\times$.
\textbf{(c)}~The two coefficients with the intervals the deposited data actually support. Bars
are $95\%$ cluster bootstrap intervals over the $19$ molecules ($B=20\,000$, cluster $=$
molecule): $\chi$ gives $0.901$ $[0.805,0.959]$, $\mathcal F_1$ gives $0.862$ $[0.709,0.929]$.
\emph{The molecular comparison does not resolve the two resources}: $\Delta\rho=0.039$,
Williams $t=0.94$ on $35$ df, $p=0.35$ (Steiger $z=0.94$), bootstrap $95\%$ CI of the difference
$[-0.070,+0.189]$; neither does the pooled comparison ($\Delta\rho=0.123$, $t=1.50$, $p=0.14$).
The $n=74$ entries are point estimates---only the difference was resampled. The dashed line is
the ceiling $|\rho|\le0.986$ imposed by the ties in $|\mathcal S|$, against which $\chi$ reaches
$91\%$ and $\mathcal F_1$ $87\%$ of what is attainable. What discriminates between the two
candidate resources is therefore not this figure but the stratified Hubbard sweep of
Fig.~\ref{fig:decoupling}(b,d).
\textbf{(d)}~The same $38$ molecular points in the rank space in which a Spearman coefficient is
actually computed, against the identity line. The $11$ floor points share midrank $6.0$---the
vertical block at the left---which is the origin of the ceiling quoted in~(c).
\emph{Provenance and caveats.} The three ``chain'' rows are the same states as the $U=8$
half-filled entries of the Hubbard sweep (identical $|\mathcal S|$ and $\chi$, identical
$\mathcal F_1$ to the precision at which they are stored), re-plotted as the ladder's matched
partner; the $74$ rows are therefore $71$ distinct states and three of them enter the pooled
coefficients twice. Recomputed over the $71$ distinct states the pooled coefficients are
$\rho(\chi,|\mathcal S|)=+0.746$ and $\rho(\mathcal F_1,|\mathcal S|)=+0.592$. The $38$
molecular points are $19$ molecules $\times$ $2$ geometries, i.e.\ dependent pairs (effective
$n\approx31.3$), which is why every molecular interval here is clustered by molecule.}
\end{figure*}
